\section{Końce i kokońce}\label{cat:sec:ends-coends}
Końce i kokońce pozwalają opisać konstrukcje uniwersalne zależne jednocześnie kowariantnie i kontrawariantnie od tego samego argumentu. W szczególności zbiór transformacji naturalnych jest końcem, a iloraz występujący w punktowym lewym rozszerzeniu Kana jest kokońcem. Rozdział rozwija pojęcia z rozdziałów~\ref{sec:naturalne}, \ref{sec:granice} i~\ref{cat:sec:kan-extensions}. Pracujemy z małą kategorią $\mathcal A$ oraz lokalnie małą kategorią $\mathcal C$, a wybory obiektów uniwersalnych rozpatrujemy w ustalonym większym uniwersum, jak w rozdziale~\ref{cat:sec:kan-extensions}.

\subsection{Dinaturalność}
Niech $T:\mathcal A^{\op}\times\mathcal A\to\mathcal C$ będzie funktorem dwóch zmiennych. Dla $f:a\to b$ i $g:c\to d$ piszemy $T(f,g):T(b,c)\to T(a,d)$, pomijając symbol $\op$ w pierwszej współrzędnej. W szczególności
\[
 T(\id_a,f):T(a,a)\to T(a,b),\qquad
 T(f,\id_b):T(b,b)\to T(a,b).
\]
Pierwsza zmienna jest kontrawariantna, druga kowariantna. Z funktorialności wynika przemienność działań w różnych zmiennych: obie drogi od $T(b,c)$ do $T(a,d)$, stosujące kolejno $f$ i $g$, są równe $T(f,g)$.

\begin{definicja}[Transformacja dinaturalna]\label{cat:def:dinatural}
Dla funktorów $S,T:\mathcal A^{\op}\times\mathcal A\to\mathcal C$ transformacją dinaturalną $\alpha:S\dashrightarrow T$ nazywamy rodzinę morfizmów $\alpha_a:S(a,a)\to T(a,a)$ spełniającą dla każdego $f:a\to b$ równanie
\begin{equation}\label{cat:eq:dinatural}
 T(\id_a,f)\alpha_a S(f,\id_a)
 =T(f,\id_b)\alpha_b S(\id_b,f).
\end{equation}
Obie strony są morfizmami $S(b,a)\to T(a,b)$.
\end{definicja}
\[
\begin{array}{ccccc}
 && S(b,a) &&\\
 &\swarrow &&\searrow &\\
 S(a,a)&&&&S(b,b)\\
 \downarrow\alpha_a&&&&\downarrow\alpha_b\\
 T(a,a)&&&&T(b,b)\\
 &\searrow&&\swarrow&\\
 &&T(a,b)&&
\end{array}
\]
Górne strzałki są odpowiednio $S(f,\id_a)$ i $S(\id_b,f)$, dolne odpowiednio $T(\id_a,f)$ i $T(f,\id_b)$. Nie zakładamy istnienia funktora diagonalnego $a\mapsto T(a,a)$: morfizm $f:a\to b$ nie daje na ogół morfizmu $T(a,a)\to T(b,b)$.

\begin{stwierdzenie}\label{cat:prop:dinatural-natural}
Ograniczenie transformacji naturalnej $\beta:S\Rightarrow T$ do par $(a,a)$ jest transformacją dinaturalną.
\end{stwierdzenie}
\begin{proof}
Naturalność względem morfizmów $(f,\id_a)$ oraz $(\id_b,f)$ daje
\[
 \beta_{a,a}S(f,\id_a)=T(f,\id_a)\beta_{b,a},
 \qquad
 \beta_{b,b}S(\id_b,f)=T(\id_b,f)\beta_{b,a}.
\]
Po złożeniu z odpowiednimi dolnymi strzałkami obie strony~\eqref{cat:eq:dinatural} stają się $T(f,f)\beta_{b,a}$ przez funktorialność $T$.
\end{proof}
\begin{uwaga}\label{cat:rem:dinatural-composition}
Dowolne transformacje dinaturalne nie muszą składać się do transformacji dinaturalnej. Można natomiast składać transformację dinaturalną z ograniczeniem transformacji naturalnej, po obu stronach.
\end{uwaga}
Istotnie, dla $\alpha:S\dashrightarrow T$ i naturalnej $\beta:T\Rightarrow U$ mamy
\[
 U(\id_a,f)\beta_{a,a}\alpha_a S(f,\id_a)
 =\beta_{a,b}T(\id_a,f)\alpha_a S(f,\id_a)
 =\beta_{a,b}T(f,\id_b)\alpha_b S(\id_b,f),
\]
a ostatnia strona jest równa $U(f,\id_b)\beta_{b,b}\alpha_b S(\id_b,f)$ przez naturalność $\beta$. Dla naturalnej $\gamma:R\Rightarrow S$ stosujemy analogicznie równania $\gamma_{a,a}R(f,\id_a)=S(f,\id_a)\gamma_{b,a}$ i $\gamma_{b,b}R(\id_b,f)=S(\id_b,f)\gamma_{b,a}$, po czym używamy~\eqref{cat:eq:dinatural}.
\begin{przyklad}[Brak domknięcia na składanie]\label{cat:ex:dinatural-noncomposition}
Niech $\mathcal A$ będzie kategorią $0\to1$ z nietrywialną strzałką $f$, a $\mathcal C=\Set$. Niech $S,T,U$ będą funktorami na $\mathcal A^{\op}\times\mathcal A$. Bierzemy $S$ stały o wartości $\{*\}$, $T(1,0)=\varnothing$ i pozostałe trzy wartości $T$ singletonowe. Niech $U(1,0)=\varnothing$, $U(0,0)=U(1,1)=\{*\}$ i $U(0,1)=\{0,1\}$. Dwie strzałki z diagonalnych wartości $U$ do $U(0,1)$ wysyłają $*$ odpowiednio na $0$ i $1$. Pozostałe nietrywialne strzałki o dziedzinie pustej są jedynymi funkcjami z pustego zbioru. Są to funktory, bo jedyne dwie nietrywialne drogi długości dwa mają pustą dziedzinę.
Rodziny jedynych map singletonów $\alpha_a:S(a,a)\to T(a,a)$ i $\beta_a:T(a,a)\to U(a,a)$ są dinaturalne: dla $\alpha$ obie strony równania mają wspólną singletonową przeciwdziedzinę, dla $\beta$ wspólną pustą dziedzinę. Jednak dla $\beta_a\alpha_a$ obie strony równania dla $f$ są mapami $\{*\}\to\{0,1\}$ o różnych wartościach. Zatem rodzina złożeń nie jest dinaturalna.
\end{przyklad}

\subsection{Własności uniwersalne końca i kokońca}
\begin{definicja}[Stożki i kostożki dinaturalne]\label{cat:def:wedges}
Stożkiem dinaturalnym nad $T$ o wierzchołku $X$ nazywamy rodzinę $p_a:X\to T(a,a)$ spełniającą
\begin{equation}\label{cat:eq:wedge}
 T(\id_a,f)p_a=T(f,\id_b)p_b\qquad(f:a\to b).
\end{equation}
Kostożek dinaturalny o wierzchołku $Y$ jest rodziną $i_a:T(a,a)\to Y$ spełniającą
\begin{equation}\label{cat:eq:cowedge}
 i_a T(f,\id_a)=i_b T(\id_b,f)\qquad(f:a\to b).
\end{equation}
Są to transformacje dinaturalne odpowiednio ze stałego funktora do $T$ i z $T$ do stałego funktora.
\end{definicja}
\begin{definicja}[Koniec i kokoniec]\label{cat:def:end-coend}
Końcem funktora $T$ nazywamy obiekt $E$ z dinaturalnym stożkiem $p_a:E\to T(a,a)$, taki że dla każdego stożka dinaturalnego $q_a:X\to T(a,a)$ istnieje dokładnie jeden $h:X\to E$ spełniający $p_a h=q_a$ dla wszystkich $a$. Piszemy
\[
 E=\int_{a\in\mathcal A}T(a,a).
\]
Kokońcem nazywamy obiekt $Q$ z dinaturalnym kostożkiem $i_a:T(a,a)\to Q$, taki że dla każdego kostożka dinaturalnego $j_a:T(a,a)\to Y$ istnieje dokładnie jeden $h:Q\to Y$ spełniający $h i_a=j_a$ dla wszystkich $a$. Piszemy
\[
 Q=\int^{a\in\mathcal A}T(a,a).
\]
Symbol całki jest tu oznaczeniem konstrukcji kategoryjnej, nie całki względem miary.
\end{definicja}
\begin{stwierdzenie}[Jednoznaczność]\label{cat:prop:end-coend-unique}
Jeżeli koniec (odpowiednio kokoniec) istnieje, jest jednoznaczny z dokładnością do jedynego izomorfizmu zgodnego ze stożkami (odpowiednio kostożkami) uniwersalnymi.
\end{stwierdzenie}
\begin{proof}
Dla dwóch końców $(E,p),(E',p')$ własności uniwersalne wyznaczają $h:E\to E'$ i $k:E'\to E$ z $p'_a h=p_a$, $p_a k=p'_a$. Mamy $p_a kh=p_a$ i $p'_a hk=p'_a$ dla każdego $a$, więc jednoznaczność daje $kh=\id_E$, $hk=\id_{E'}$. Dla kokońców $(Q,i),(Q',i')$ otrzymujemy $h:Q\to Q'$, $k:Q'\to Q$ z $h i_a=i'_a$, $k i'_a=i_a$. Równania $kh i_a=i_a$ i $hk i'_a=i'_a$ dają analogicznie wzajemną odwrotność. Jednoznaczność samego porównania jest częścią własności uniwersalnej.
\end{proof}
\begin{stwierdzenie}[Funktorialność]\label{cat:prop:end-coend-functorial}
Jeśli końce wszystkich rozważanych funktorów istnieją, wybór końców wyznacza funktor na kategorii tych funktorów i transformacji naturalnych. Analogicznie dla kokońców.
\end{stwierdzenie}
\begin{proof}
Niech $\beta:T\Rightarrow U$ będzie transformacją naturalną, a $p^T,p^U$ stożkami końców. Rodzina $\beta_{a,a}p^T_a$ jest dinaturalna na mocy uwagi~\ref{cat:rem:dinatural-composition}. Istnieje zatem jedyny $E(\beta):E(T)\to E(U)$ z
\[
 p^U_a E(\beta)=\beta_{a,a}p^T_a.
\]
Dla $\beta=\id$ równanie spełnia identyczność. Dla $\gamma:U\Rightarrow V$ złożenie $E(\gamma)E(\beta)$ ma po postkompozycji z $p^V_a$ wartość $\gamma_{a,a}\beta_{a,a}p^T_a$, więc jest równe $E(\gamma\beta)$. Dla kokońców rodzina $i^U_a\beta_{a,a}$ jest dinaturalnym kostożkiem. Jedyny $Q(\beta):Q(T)\to Q(U)$ spełnia $Q(\beta)i^T_a=i^U_a\beta_{a,a}$. Identyczności i składanie zachowują się z tego samego argumentu: oba złożenia z $i^T_a$ dają $i^V_a\gamma_{a,a}\beta_{a,a}$.
\end{proof}

\subsection{Konstrukcja przez equalizatory i koequalizatory}
\begin{twierdzenie}[Istnienie końca]\label{cat:thm:end-equalizer}
Załóżmy, że w $\mathcal C$ istnieją produkty
\[
 P=\prod_{a\in\Ob(\mathcal A)}T(a,a),
 \qquad D=\prod_{f:a\to b}T(a,b).
\]
Oznaczmy projekcje przez $\pi_a$ i $\pi_f$. Niech $s,t:P\to D$ będą wyznaczone przez
\[
 \pi_f s=T(\id_a,f)\pi_a,\qquad
 \pi_f t=T(f,\id_b)\pi_b .
\]
Jeżeli istnieje equalizator $e:E\to P$ pary $s,t$, to $E$ wraz z $p_a=\pi_a e$ jest końcem $T$. W szczególności końce istnieją w każdej kategorii zupełnej.
\end{twierdzenie}
\begin{proof}
Morfizmy $s,t$ istnieją i są jedyne na mocy własności uniwersalnej produktu $D$. Równość $se=te$ po postkompozycji z $\pi_f$ daje~\eqref{cat:eq:wedge} dla $p_a=\pi_a e$. Dla dowolnego stożka dinaturalnego $q_a:X\to T(a,a)$ istnieje jedyny $q:X\to P$ z $\pi_a q=q_a$. Dla każdego $f:a\to b$ mamy
\[
 \pi_f s q=T(\id_a,f)q_a=T(f,\id_b)q_b=\pi_f t q .
\]
Jednoznaczność w produkcie $D$ daje $sq=tq$. Z własności equalizatora istnieje jedyny $h:X\to E$ z $eh=q$; wtedy $p_a h=q_a$. Jeśli $h'$ również spełnia $p_a h'=q_a$, to $\pi_a eh'=q_a$ dla każdego $a$, stąd $eh'=q$, a jednoznaczność w equalizatorze daje $h'=h$. Małość $\mathcal A$ gwarantuje, że oba produkty są małe; zupełność daje je oraz equalizator.
\end{proof}
\begin{twierdzenie}[Istnienie kokońca]\label{cat:thm:coend-coequalizer}
Załóżmy, że istnieją koprodukty
\[
 P'=\coprod_{a\in\Ob(\mathcal A)}T(a,a),
 \qquad D'=\coprod_{f:a\to b}T(b,a).
\]
Niech $j_a,j_f$ oznaczają ich morfizmy kanoniczne. Definiujemy $s',t':D'\to P'$ przez
\[
 s'j_f=j_a T(f,\id_a),\qquad
 t'j_f=j_b T(\id_b,f).
\]
Jeśli istnieje koequalizator $q:P'\to Q$ pary $s',t'$, to $Q$ wraz z $i_a=qj_a$ jest kokońcem $T$. W szczególności kokońce istnieją w każdej kategorii kozupełnej.
\end{twierdzenie}
\begin{proof}
Istnienie i jednoznaczność $s',t'$ wynika z własności uniwersalnej koproduktu $D'$. Równość $qs'=qt'$ po prekompozycji z $j_f$ daje~\eqref{cat:eq:cowedge}. Dla dowolnego kostożka dinaturalnego $k_a:T(a,a)\to Y$ istnieje jedyny $k:P'\to Y$ z $kj_a=k_a$. Mamy
\[
 ks'j_f=k_aT(f,\id_a)=k_bT(\id_b,f)=kt'j_f .
\]
Z jednoznaczności w koprodukcie $D'$ wynika $ks'=kt'$. Własność koequalizatora daje jedyny $h:Q\to Y$ z $hq=k$, więc $hi_a=k_a$. Jeśli $h'$ spełnia $h'i_a=k_a$, to $h'qj_a=k_a$ dla wszystkich $a$, a stąd $h'q=k$. Jednoznaczność w koequalizatorze daje $h'=h$. Oba koprodukty są małe, więc kozupełność zapewnia istnienie wszystkich potrzebnych obiektów.
\end{proof}
\begin{wniosek}[Modele zbiorowe]\label{cat:cor:end-coend-set}
Dla $T:\mathcal A^{\op}\times\mathcal A\to\Set$ istnieją koniec i kokoniec. Można je opisać jako
\[
 \int_a T(a,a)=\left\{(x_a)_a\in\prod_a T(a,a):
 T(\id_a,f)(x_a)=T(f,\id_b)(x_b)\ \forall f:a\to b\right\}
\]
oraz
\[
 \int^a T(a,a)=\left(\coprod_a T(a,a)\right)/{\sim},
\]
gdzie $\sim$ jest najmniejszą relacją równoważności zawierającą dla każdego $f:a\to b$ i $z\in T(b,a)$ parę
\[
 (a,T(f,\id_a)(z))\sim(b,T(\id_b,f)(z)).
\]
\end{wniosek}
\begin{proof}
W $\Set$ equalizator pary funkcji jest podzbiorem punktów, na których wartości funkcji są równe (definicja~\ref{def:equalizer}). Zastosowanie tego do $s,t$ z twierdzenia~\ref{cat:thm:end-equalizer} daje pierwszy wzór. Koequalizator pary funkcji jest ilorazem przeciwdziedziny przez najmniejszą relację równoważności zawierającą pary ich wartości w każdym punkcie dziedziny (definicja~\ref{def:coeq}). Dla $s',t'$ z twierdzenia~\ref{cat:thm:coend-coequalizer} otrzymujemy dokładnie wskazane pary. Morfizmy uniwersalne są odpowiednio projekcjami ograniczonymi do zbioru zgodnych rodzin i funkcjami $x\mapsto[a,x]$.
\end{proof}
\begin{przyklad}[Kategorie dyskretne i puste]\label{cat:ex:end-discrete}
Jeśli $\mathcal A$ jest dyskretna, równania~\eqref{cat:eq:wedge} i~\eqref{cat:eq:cowedge} dla identyczności są tautologiami. Koniec jest wówczas produktem wszystkich $T(a,a)$, a kokoniec koproduktem. Dla pustej kategorii oznacza to odpowiednio obiekt końcowy i początkowy, o ile istnieją. Dla kategorii z jednym obiektem i jego identycznością obie konstrukcje są kanonicznie izomorficzne z jedyną wartością $T(a,a)$.
\end{przyklad}
\begin{przyklad}[Funktor zależny od jednej zmiennej]\label{cat:ex:end-one-variable}
Jeśli $T(a,b)=F(b)$ dla $F:\mathcal A\to\mathcal C$ (również na morfizmach $T(f,g)=F(g)$), to~\eqref{cat:eq:wedge} przyjmuje postać $F(f)p_a=p_b$, a~\eqref{cat:eq:cowedge} postać $i_a=i_b F(f)$. Są to zwykłe warunki stożka i kostożka nad $F$. Zatem koniec i kokoniec tego $T$ są odpowiednio granicą i kogranicą $F$, również pod względem istnienia. Dla $T(a,b)=G(a)$ z $G:\mathcal A^{\op}\to\mathcal C$ otrzymujemy granicę i kogranicę $G$ w kategorii przeciwnej indeksów: warunki to $p_a=G(f)p_b$ oraz $i_a G(f)=i_b$.
\end{przyklad}

\subsection{Transformacje naturalne jako końce}
\begin{twierdzenie}\label{cat:thm:nat-as-end}
Dla funktorów $F,G:\mathcal A\to\mathcal C$ istnieje bijekcja
\[
 \Nat(F,G)\cong\int_{a\in\mathcal A}\Hom_{\mathcal C}(F(a),G(a)).
\]
Po prawej stronie bierzemy koniec w $\Set$ funktora $T(a,b)=\Hom_{\mathcal C}(F(a),G(b))$. Bijekcja jest naturalna kontrawariantnie względem $F$ i kowariantnie względem $G$.
\end{twierdzenie}
\begin{proof}
Definiujemy $T(f,g)(h)=G(g)hF(f)$ dla $f:a\to a'$, $g:b\to b'$ i $h:F(a')\to G(b)$. Identyczności zachowują się. Dla $f':a'\to a''$, $g':b'\to b''$ i $h:F(a'')\to G(b)$ mamy
\[
 T(f,g')\bigl(T(f',g)(h)\bigr)
 =G(g')G(g)hF(f')F(f)
 =T(f'f,g'g)(h),
\]
gdzie w pierwszej operacji $T(f',g)$ i w drugiej $T(f,g')$ współrzędne mają zgodne typy. Wynika to z funktorialności $F,G$ oraz łączności składania. Zatem $T$ jest funktorem. Z~\ref{cat:cor:end-coend-set} element końca jest rodziną $\alpha_a:F(a)\to G(a)$ spełniającą dla każdego $f:a\to b$
\[
 G(f)\alpha_a=\alpha_bF(f).
\]
Jest to dokładnie warunek z definicji~\ref{def:naturalna}; obie strony bijekcji są więc tym samym zbiorem rodzin. Dla $\mu:F'\Rightarrow F$ i $\nu:G\Rightarrow G'$ transformacja naturalna odpowiednich funktorów dwóch zmiennych ma składowe $h\mapsto\nu_b h\mu_a$. Naturalność względem $(f,g)$ wynika z równań $F(f)\mu_a=\mu_{a'}F'(f)$ oraz $G'(g)\nu_b=\nu_{b'}G(g)$. Indukowana funkcja końców wysyła rodzinę $(\alpha_a)$ na $(\nu_a\alpha_a\mu_a)$ na mocy stwierdzenia~\ref{cat:prop:end-coend-functorial}; jest to dokładnie składanie transformacji naturalnych.
\end{proof}
\begin{uwaga}
W przypadku funktorów do kategorii $\mathcal C$ nieposiadającej dodatkowej struktury powyższy koniec jest obiektem $\Set$, a nie obiektem $\mathcal C$. Małość $\mathcal A$ i lokalna małość $\mathcal C$ zapewniają, że $\Nat(F,G)$ jest zbiorem. Wzbogaconą wersję tej konstrukcji omówimy oddzielnie.
\end{uwaga}
